\centerline{\bf \Huge Банки. Жесть}
\centerline{\bf \Huge }

\begin{flushright}
        \scriptsize{\textbf{ww}}
\end{flushright}

\problem{\textbf{1}} 15-ого апреля планируется взять кредит в банке на 700 тысяч рублей на ($n$ + 1) месяц.
Условия его возврата таковы:

— 1-го числа каждого месяца долг увеличивается на 1\% по сравнению с концом предыдущего
месяца;

— со 2-го по 14-е число каждого месяца необходимо выплатить одним платежом часть долга;

— 15-го числа каждого месяца с первого по $n$-ый долг должен быть на одну и ту же сумму
меньше долга на 15-е число предыдущего месяца;

— 15-го числа $n$-го месяца долг составлял 300 тысяч рублей;

— к 15-му числу ($n$ + 1)-го месяца долг должен быть погашен полностью.

Найдите $n$, если банку всего было выплачено 755 тысяч рублей.

\problem{\textbf{2}} В июле планируется взять кредит в банке на сумму 6 млн рублей на срок 15 лет. Условия его возврата таковы:

— каждый январь долг возрастает на $x$\% по сравнению с концом предыдущего года;

— с февраля по июнь каждого года необходимо выплатить часть долга;

— в июле каждого года долг должен быть на одну и ту же величину меньше долга на июль
предыдущего года.

Найдите $x$, если известно, что наибольший платёж по кредиту составит не более 1,9 млн
рублей, а наименьший — не менее 0,5 млн рублей.


\problem{\textbf{3}} В июле 2016 года планируется взять кредит на пять лет на сумму 1100 тыс. рублей.
Условия его возврата таковы:

— каждый январь долг будет возрастать на 20\% по сравнению с концом предыдущего года;

— с февраля по июнь каждого года необходимо выплатить одним платежом часть долга;

— в 2017, 2018 и 2019 годах сумма долга не изменяется;

— платежи в 2020 и в 2021 годах должны быть равны;

— к июлю 2021 года долг должен быть выплачен полностью.
На сколько рублей будут отличаться первый и последний платежи?

\problem{\textbf{4}} На последние два года обучения в университете студент взял образовательный кредит.
Условия пользования кредитом таковы:

— в сентябре каждого года в течение обучения студента банк перечисляет на счет университета сумму, равную стоимости годового обучения в университете;

— один раз в ноябре каждого года пользования кредитом банк начисляет 20\% на текущий
долг клиента;

— каждый год в течение обучения в декабре студент вносит некоторую неизменную сумму
в счет погашения кредита;

— после окончания обучения в течение еще двух лет банк продолжает в ноябре каждого
года начислять 20\% на оставшуюся сумму долга, но теперь студент обязан выплачивать кредит
равными платежами, в пять раз превышающими платеж во время обучения.

Сколько рублей составит переплата по такому кредиту, если год обучения в университете
стоит 402500 рублей?

\problem{\textbf{5}} В июле 2025 года планируется взять кредит на $n$ лет на неизвестную сумму. Условия его возврата таковы:

— в январе каждого года, начиная с 2026, долг возрастает на 10\% по сравнению с концом
предыдущего года;

— в июле 2026, 2027, 2028 и 2029 годов долг должен быть на 20 тыс. рублей меньше долга
на июль предыдущего года;

— в июле каждого года, начиная с 2030, долг должен быть на 10 тыс. рублей меньше долга
на июль предыдущего года;

— к июлю (2025 + $n$)-го года долг должен быть полностью погашен.

На какое наименьшее количество лет должен быть взять кредит, чтобы седьмой платеж был
не менее 30 тыс. рублей.


